\section{Where the record meets the owner's other workbenches}\label{sec:overlaps}

The audit note \note{47\_} records every crossing between the owner's workbenches that the audit found, as typed edges:
\begin{itemize}
\item \emph{P}: a proof-level import;
\item \emph{I}: the same identity or object on both sides;
\item \emph{A}: an analogy or a shared number only.
\end{itemize}
The table lists the edges that touch the $S^6$ material, and adds one found for this reader (marked ``new'').

\begin{center}\small
\begin{longtable}{@{}>{\raggedright\arraybackslash}p{0.13\linewidth}>{\raggedright\arraybackslash}p{0.44\linewidth}>{\raggedright\arraybackslash}p{0.07\linewidth}>{\raggedright\arraybackslash}p{0.28\linewidth}@{}}
\toprule
Edge & What crosses & Type & Status\\
\midrule
\endhead
$S^6\to$ YM & The imaginary part of the period matrix, used as a magnetic background in the Yang--Mills workbench. Only $D_{\mathrm{YM}}=L_Tq_T+6m_T^2=-\det_{\R}\Pi$ enters (Section~\ref{sec:candidate}) & P & audited (audit note \note{28\_}, Lemma~28.2; \cite{YMReader}, §§6.3 and~8)\\
$S^6\to$ ES & The integral monodromies $A_1,A_2,A_\infty$ of the $(3,4,\infty)$ period system, with $A_1^3=A_2^4=A_1A_2A_\infty=I$, in the Erdős--Straus package ES-09. Nothing else from the construction is imported & P & audited for odd levels, which the application uses; false at even levels (\cite{ESReader}, §5.2)\\
ES $\leftrightarrow$ $S^6$ & The Niemeier lattice with root system $A_5^4D_4$ and glue index $72$: the Erdős--Straus result R10, and S6-4 above & I & the Erdős--Straus glue code verified; the $S^6$ code is in the record's archive \fn{28\_}, and the two were not compared\\
$S^6$ and the Jacobian map $\to$ NS & Fluid coordinates, and gauge-scalar equations from the $S^6$ material carried to Navier--Stokes on $\mathbb T^3$, inside the zeta repository & P (displayed calculations) & located (audit note \note{47\_}, edge 13)\\
$S^6\leftrightarrow$ ES (new) & The number $5184=72^2$. It is $\det I$ in S6-5, where $I$ has Gram matrix $6G_{D_4}$, so $\det I=6^4\det D_4$; and it is $|\operatorname{disc}(A_5^4D_4)|=(\det A_5)^4\det D_4$ in S6-4 and R10 & A & both products $6^4\cdot4$ checked; no map between the two lattices was found\\
\bottomrule
\end{longtable}
\end{center}

Three further points of contact are not mathematical edges:
\begin{itemize}
\item \emph{Hosting.} The Yang--Mills repository has an \texttt{s6/} topic entry, which holds the key-advances reader, and the record's frozen edition of 9 September was released together with companion Yang--Mills and Navier--Stokes editions (guide \fn{29\_}).
\item \emph{The name ``Fable''.} The name labels three different things across the repositories: the three-dimensional Keller map of the Jacobian counterexample, the $S^6$ construction (``with Fable'' in older $S^6$ files), and a conductor cover of the Erdős--Straus crosswalk (\cite{ESReader}, §7). The Yang--Mills repository's attribution file, which its author designates as the clarification, says that the $S^6$ construction is a different object from the Jacobian counterexample, and that it was circulated by Levent Alpöge and produced with Claude.
\item \emph{The correction note and the zeta programme.} No crossing was found. The derived criterion (Theorem~\ref{thm:derived}) and the local algebra of Section~\ref{sec:correction} use nothing from the other workbenches.
\end{itemize}
